% From docs/: make images/daq_diagram.pdf
\documentclass[tikz,border=8mm]{standalone}
\usepackage[T1]{fontenc}
\usepackage{lmodern}
\usepackage[american]{circuitikz}
\usetikzlibrary{arrows.meta,calc}
\standaloneenv{circuitikz}
\renewcommand{\familydefault}{\sfdefault}
\ctikzset{tripoles/plain amp/height=1.0,tripoles/fd op amp/height=1.0}
\colorlet{ink}{black}
\definecolor{pcb}{HTML}{E5E5E5}
\definecolor{silicon}{HTML}{FFFFFF}
\tikzset{box/.style={draw,fill=white,minimum height=1.1cm,align=center,font=\small},
 flow/.style={-{Latex[length=2mm]},line width=.7pt},
 clock/.style={line width=.7pt,dashed},
 note/.style={align=left,font=\small},tinylabel/.style={font=\footnotesize,align=center},
 % Native differential amplifier shape, adapted to one input; P/N label the outputs.
 lvdsdriver/.style={fd op amp, circuitikz/tripoles/fd op amp/input height=0,
 circuitikz/amplifiers/plus={},circuitikz/amplifiers/minus={},fill=white},
 pisoreg/.style={muxdemux,muxdemux def={NL=1,NR=1,NT=2,NB=0,Lh=3.5,Rh=1.6,w=3.5},fill=white},
 bank/.style={draw,minimum width=4.4cm,minimum height=1.7cm,fill=white}}
\begin{document}
\begin{circuitikz}[x=1cm,y=1cm,line width=.7pt]
\path[use as bounding box] (0,2.2) rectangle (42.3,21.6);
% Two boards, with the ASIC on the left and FPGA on the right.
\draw[black!55,fill=pcb] (0,4.0) rectangle (9,17.8);
\node[anchor=north west,font=\bfseries] at (.3,17.45) {FRIDA single-chip card};
\draw[black!55,fill=silicon] (.4,4.35) rectangle (8.6,16.8);
\node[anchor=north west,font=\bfseries] at (.7,16.5) {FRIDA-1 ASIC};
\draw[black!55,fill=pcb] (13,2.7) rectangle (37.1,21.3);
\node[anchor=north west,font=\bfseries] at (13.3,21) {BDAQ53 carrier / Mercury KX1 FPGA module};
\draw[black!55,fill=silicon] (13.35,3.1) rectangle (33.65,18.1);
\node[anchor=north west,tinylabel] at (13.6,17.85) {XC7K160T-1};
\node[tinylabel] at (10.5,9.4) {\textbf{DisplayPort cable}};
% External programmable oscillator sits in the board area to the right of the FPGA.
\node[box,minimum width=2cm,minimum height=1.6cm] (si570) at (35.4,17.45) {};
\node[vcoshape,circuitikz/bipoles/vco/width=.6,fill=white] (xo) at (35.4,17.3) {};
\node[font=\ttfamily\small,fill=white,inner sep=1pt] at (35.4,17.95) {si570};
% Functional PLL internals; clock buffers are omitted from this view.
\draw[black!55] (22,15.55) rectangle (31.6,18.05);
\node[tinylabel,anchor=north west] at (22.15,17.95) {\texttt{plle2\_adv\_seq}};
\node[box,minimum width=1cm,minimum height=.55cm] (mult) at (26.9,17.45) {$\times 8$};
\node[vcoshape,circuitikz/bipoles/vco/width=.65,fill=white] (pll) at (26.9,16.3) {};
\draw[clock] (si570.west) -- (mult.east);
\draw[clock] (mult.south) -- (pll.north);
\node[tinylabel] at (35.4,19.3) {100--200 MHz\\$\approx0.09$ ppb tuning steps};
\node[box,minimum width=1.2cm,minimum height=.7cm] (serclk) at (23.2,16.3) {$\div N$};
\node[box,minimum width=1.2cm,minimum height=.7cm] (wordclk) at (30.7,16.3) {$\div(4N)$};
\draw[clock] (pll.west) -- (serclk.east);
\draw[clock] (pll.east) -- (wordclk.west);
\node[tinylabel] at (29.2,16.95) {$N=2\ldots20$};
\node[tinylabel] at (27,18.65) {100 / 320 / 960 / 1600 MBd $\to N=20 / 10 / 3 / 2$};
% Sequencer: one register outline with a top clock input; RAM-backed implementation.
\node[bank] (seq) at (30.1,12.3) {};
\node[tinylabel] at (30.1,12.3) {\texttt{seq\_gen}\\cyclic 64-bit words};
\node[tinylabel] at (30.1,10.8) {RAM-backed sequencer\\bytes 0--3 $\to$ controls};
\draw (30.45,13.15) -- (30.7,12.95) -- (30.95,13.15);
\coordinate (seqclk) at (30.7,13.15);
% Serializer: one 32-bit bus represents four parallel eight-bit lanes.
% Four physical OSERDES lanes, represented by stacked mux-shaped serializers.
\node[pisoreg,muxdemux def={NL=0,NR=0,NT=2},draw only top pins={},xscale=-1] (ser3) at (24.4,12.9) {};
\foreach \x/\y/\n in {24.2/12.7/2,24/12.5/1}{
 \node[pisoreg,muxdemux def={NL=0,NR=0,NT=0},xscale=-1] (ser\n) at (\x,\y) {};
}
\node[pisoreg,muxdemux def={NT=0},xscale=-1] (ser) at (23.8,12.3) {};
\node[font=\ttfamily\footnotesize] at (23.8,12.3) {oserdese2};
\node[tinylabel] at (23.8,10.55) {4 $\times$ 8:1 DDR; d1 first};
\draw[flow] (seq.west) -- ($(ser.lpin 1)+(.35,0)$);
\draw ($(ser.lpin 1)+(.35,0)$) -- (ser.lpin 1);
\draw (26.35,12.08) -- (26.6,12.52);
\node[tinylabel,anchor=west] at (26.63,12.55) {32};
\node[tinylabel] at (26.8,13.45) {4 $\times$ 8-bit buses};
\node[tinylabel] at (26.7,11) {Board-polarity\\correction};
% Clock wedges interrupt the actual upper edge of the visible front serializer.
\coordinate (serbitl) at ($(ser.top left)!0.30!(ser.top right)$);
\coordinate (serbitr) at ($(ser.top left)!0.42!(ser.top right)$);
\coordinate (serwordl) at ($(ser.top left)!0.61!(ser.top right)$);
\coordinate (serwordr) at ($(ser.top left)!0.73!(ser.top right)$);
\coordinate (serbitpin) at ($(serbitl)!.5!(serbitr)$);
\coordinate (serwordpin) at ($(serwordl)!.5!(serwordr)$);
\foreach \left/\right/\pin in {serbitl/serbitr/serbitpin,serwordl/serwordr/serwordpin}{
 \draw[white,line width=1.6pt] (\left) -- (\right);
 \draw (\left) -- ($(\pin)+(0,-.19)$) -- (\right);
}
\draw[clock] (wordclk.south) -- (30.7,14.6) -- (seqclk);
\draw[clock] (30.7,14.6) -- (serwordpin |- 30.7,14.6) -- (serwordpin);
\node[tinylabel,anchor=south] at (27.2,14.75) {\texttt{seq\_clk}: 10--200 MHz};
\draw[clock] (serclk.south) -- (23.2,15.05) -| (serbitpin);
\node[tinylabel,anchor=east,align=right] at (22.95,14.8) {\texttt{ser\_clk}\\40--800 MHz};
\draw[clock] (30.7,14.6) -- (33.05,14.6) -- (33.05,7.8) -- (30.1,7.8) -- (30.1,7.15);
\fill (30.7,14.6) circle (1.5pt);
% Outgoing LVDS: actual single-ended input / differential output topology.
% Four-lane stack: 0.10 cm right and 0.22 cm up per lane.
\foreach \d in {.30,.20,.10}{
 % Body silhouettes omit unused rear-layer terminal stubs.
 \draw[line width=1.4pt,fill=white] (16+\d-.833,12.3+2.2*\d) --
  (16+\d+.833,12.3+2.2*\d+.7) -- (16+\d+.833,12.3+2.2*\d-.7) -- cycle;
}
\node[lvdsdriver,xscale=-1] (tx) at (16,12.3) {};
\node[tinylabel] at (16,14.5) {4 $\times$ \texttt{obufds}\\output I/O bank: HR 15\\lvds\_25; no odelaye2};
\draw[flow] (ser.rpin 1) -- (tx.in up);
\node[tinylabel] at (19.45,13.35) {4 serial lanes};
\node[tinylabel] at (19.45,11.55) {0.625 ns / symbol};
\draw (19.45,12.08) -- (19.75,12.52);
\node[tinylabel,anchor=west] at (19.78,12.55) {4};
\foreach \d in {.30,.20,.10}{
 \draw[line width=1.4pt,fill=white] (4.5+\d-.833,12.3+2.2*\d) --
  (4.5+\d+.833,12.3+2.2*\d+.7) -- (4.5+\d+.833,12.3+2.2*\d-.7) -- cycle;
}
\node[plain amp,xscale=-1,fill=white] (rx) at (4.5,12.3) {};
\node[font=\small] at (5.05,12.55) {$+$};
\node[font=\small] at (5.05,12.05) {$-$};
\node[tinylabel] at (4.5,14.5) {\textbf{4 $\times$ LVDS RX pads}\\\texttt{seq\_init / seq\_samp}\\\texttt{seq\_comp / seq\_logic}};
\draw (tx.out up) -- (rx.in up);
\draw (tx.out down) -- (rx.in down);
\node[tinylabel,fill=white] at (10.5,13.4) {\texttt{seq\_init, seq\_samp}\\\texttt{seq\_comp, seq\_logic}};
\foreach \side in {up,down}{
 \coordinate (busmark) at (11.4,0 |- tx.out \side);
 \draw ($(busmark)+(-.12,-.18)$) -- ($(busmark)+(.12,.18)$);
 \node[tinylabel,anchor=west,inner sep=1pt] at ($(busmark)+(.18,.16)$) {4};
}

\node[font=\small] at (16.0,12.45) {$+$};
\node[font=\small] at (16.0,12.15) {$-$};
% The ADC stack is low and left so the comparator bus runs directly into its mux.
% Rear silhouettes shift up and left; the front trapezoid contains the A/D label.
\foreach \d in {.36,.24,.12,0}{
 \draw[line width=1.4pt,fill=white]
  (1.55-\d,6.3+1.5*\d-.48) -- (2.95-\d,6.3+1.5*\d-.78) --
  (2.95-\d,6.3+1.5*\d+.78) -- (1.55-\d,6.3+1.5*\d+.48) -- cycle;
}
\node[font=\footnotesize] at (2.25,6.3) {A/D};
\node[tinylabel] at (1.45,8.05) {16 ADCs};
% Four single-ended clock trees enter the top of the front ADC.
\coordinate (adcclk) at (2.25,6.93);
\draw (rx.out) -- (2.25,12.3) -- (adcclk);
\draw[white,line width=1.6pt] (2.11,6.90) -- (2.39,6.96);
\draw (2.11,6.90) -- (2.25,6.70) -- (2.39,6.96);
\draw (2.95,12.08) -- (3.25,12.52);
\node[tinylabel,anchor=west] at (3.28,12.55) {4};
\node[tinylabel,anchor=west,align=left] at (4.2,10) {4 single-ended\\clock trees};
\draw[flow] (.5,6.48) -- (1.55,6.48);
\draw[flow] (.5,6.12) -- (1.55,6.12);
\node[tinylabel] at (.85,5.6) {$V_{\mathrm{in}}$};
% Comparator bus, not the on-chip SAR output-code bus.
\node[muxdemux,muxdemux def={NL=1,NR=1,NB=1,NT=0,Lh=2.2,Rh=1.2,w=1.5},fill=white] (compmux) at (4.15,6.3) {};
\draw (2.95,6.3) -- (compmux.lpin 1);
\draw (3.13,6.08) -- (3.43,6.52);
\node[tinylabel,anchor=west] at (3.41,6.55) {16};
\node[tinylabel] at (4.15,6.3) {16:1};
\draw (compmux.bpin 1) -- (4.15,5.15);
\draw (3.98,5.32) -- (4.32,5.53);
\node[tinylabel,anchor=west] at (4.38,5.42) {4};
\node[lvdsdriver] (asic_tx) at (6.05,6.3) {};
\draw[flow] (compmux.rpin 1) -- ($(asic_tx.center)+(-.73,0)$);
\node[tinylabel] at (6.05,4.8) {\textbf{LVDS TX pad}};
% Returning comparator bit, followed by two physical delay elements.
\node[plain amp,fill=white] (fpga_rx) at (16,6.3) {};
\node[font=\small] at (15.45,6.55) {$+$};
\node[font=\small] at (15.45,6.05) {$-$};
\draw (asic_tx.out up) -- (fpga_rx.in up);
\draw (asic_tx.out down) -- (fpga_rx.in down);
\node[tinylabel] at (10.5,7.4) {\texttt{comp\_out}};
\node[tinylabel] at (16,4.85) {\texttt{ibufds}};
\node[font=\small] at (6.05,6.45) {$+$};
\node[font=\small] at (6.05,6.15) {$-$};
\node[box,minimum width=2.6cm] (id0) at (20.1,6.3) {\texttt{idelaye2}\\$\tau_0+n_0\Delta t$};
\node[box,minimum width=2.6cm] (id1) at (22.7,6.3) {\texttt{idelaye2}\\$\tau_1+n_1\Delta t$};
\draw[flow] (fpga_rx.out) -- (id0.west);
\node[bank] (fast) at (30.1,6.3) {};
\node[tinylabel] at (30.1,6.3) {\texttt{fast\_spi\_rx}\\17-bit shift register + CDC FIFO\\one bit / rising sclk edge};
\draw (29.85,7.15) -- (30.1,6.95) -- (30.35,7.15);
\node[tinylabel,anchor=east] at (30,7.5) {\texttt{sclk}};
\node[tinylabel,anchor=west] at (28,6.82) {\texttt{sen}};
\draw[flow] (id1.east) -- node[above,tinylabel]{\texttt{sdi}} (fast.west);
\draw[flow] ($(seq.west)+(0,-.5)$) -- (27.6,11.8) -- (27.6,6.82) -- (27.9,6.82);
\node[tinylabel,fill=white,align=left] at (30.2,9.2) {\texttt{seq\_out[32]} $\to$ \texttt{sen}\\17-word capture window\\start word $w$: 5 ns steps};
% Delay control: independent fixed reference, not the variable sequencer clock.
\node[box,minimum width=7cm] (ctl) at (22,8.75) {\texttt{gpio1 + idelayctrl}\\$N=0\ldots62$, $n_0=\lfloor N/2\rfloor$, $n_1=\lceil N/2\rceil$\\refclk: fixed 200 MHz from \texttt{plle2\_base\_comm}};
\draw[flow] (20.1,8.1) -- (id0.north);
\draw[flow] (22.7,8.1) -- (id1.north);
% fast_spi_rx owns CDC; its 32-bit output feeds a separate width-conversion FIFO.
\node[box,minimum width=5.6cm] (fifo) at (30.1,4.0) {\texttt{fifo\_32\_to\_8}\\bus\_clk $\approx142.86$ MHz};
\node[box,minimum width=2.8cm] (ethernet) at (35.4,4.0) {Ethernet PHY\\board I/O};
\draw[flow] (fast.south) -- (fifo.north);
\draw[flow] (fifo.east) -- node[above,tinylabel]{SiTCP} (ethernet.west);
\draw[flow] (ethernet.east) -- node[pos=.55,above,tinylabel]{to PC via TCP} (41.8,4.0);
\end{circuitikz}
\end{document}
